\subsection{Lower bounds for additive depth}
\label{sec:pn-depth-lower}

\subsubsection{A fixed-stabilizer lower bound}

Theorem~\ref{thm:pn-relative} needs its relative floor. With a fixed
stabilizer alone, logarithmic depth can force a cost linear in the
context length, even at zero attention scores.
The following construction concerns weight-only preprocessing, as in the
input-probe model. It also applies when the unknown signs arrive as a
fresh query. It is not a lower bound against an index that has already
read and summarized these particular signs during prefill.

\begin{theorem}[Additive pre-norm blocks can reveal a majority]
\label{thm:pn-absolute-lower}
For every $T\ge2$, there is a width-three instance of
\eqref{eq:pn-att}--\eqref{eq:pn-ff}, with $s=4$, stabilizer
$\epsilon=3$ at every normalization, zero attention scores, and
$O(\log T)$ blocks, whose bounded two-logit exact-token query complexity
is at least $T/2^{17}$ in the worst context. All nonzero weights are in
$\{1,2,-2\}$ except a dyadic final readout coefficient. The input alphabet
is $\{(-1,1,-1),(1,1,-1)\}$, and every normalization denominator is at
least $\sqrt3$.
\end{theorem}
\paragraph{Proof idea.}
Uniform attention writes a hidden input average into one coordinate.
Subsequent normalized residual updates amplify small averages until their
sign becomes visible at a bounded readout. A fourth-moment estimate and the
transcript score convert this sensitivity into a linear query bound.

\begin{proof}
Write the unknown first coordinates as $x_1,\ldots,x_T$. Initially every
coordinate has magnitude one, so the first normalization denominator is
exactly two. Use zero queries and keys. The sole nonzero value row is
row two, with coefficients $2$ on coordinate one and $-2$ on coordinate
two. The attention residual therefore changes the stream at position
$t$ to
\[
 (x_t,M_t,-1),\qquad M_t=\frac1t\sum_{j\le t}x_j.
\]
Take the tanh update of this first block to be zero. All later attention
updates are zero. In each later feed-forward sublayer, the sole nonzero
row is row two, selecting its own normalized coordinate with weight one.
At the last position this produces the scalar recurrence
\[
 F_0(z)=z,\qquad
 F_{r+1}(z)=F_r(z)+
   \tanh\!\left(\frac{F_r(z)}
      {\sqrt{11/3+F_r(z)^2/3}}\right),
 \qquad z=M_T.
\]
Every recurrence is odd and increasing. If $0\le y\le1$, its normalized
argument is at least $y/2$ and at most one. The inequality
$\tanh v\ge v/2$ for $0\le v\le1$ therefore gives an increment at least
$y/4$. If $y\ge1$, the argument is at least $1/2$, so the increment is
at least $1/4$.

Set
\[
 k=\left\lceil\log_{5/4}(2\sqrt T)\right\rceil,
 \qquad d=2^{\lceil\log_2(2k)\rceil}.
\]
Thus $2k\le d<4k$. Starting at $z\ge1/(2\sqrt T)$, the recurrence
reaches one within $k$ steps. During its remaining $d-k$ steps it
increases by at least $1/4$ per step. Consequently
$F_d(z)\ge d/8$. The trivial upper bound $|F_d(z)|\le1+d\le2d$ holds
for all $|z|\le1$. Take logits $\pm F_d(M_T)/(2d)$, a legal dyadic
readout. Their output-sign mean is
\[
 g(M_T)=\tanh(F_d(M_T)/(2d)).
\]
For $z\ge1/(2\sqrt T)$, one has $g(z)\ge1/32$. Oddness gives the
corresponding negative inequality.

Let the input signs be independent with common mean $t$, and put
$H(t)=\mathbb E_tg(M_T)$. Differentiating the finite input sum at zero
gives
\[
 H'(0)=T\mathbb E_0[M_Tg(M_T)].
\]
Now $\mathbb E_0M_T^2=1/T$ and
$\mathbb E_0M_T^4=(3T^2-2T)/T^4\le3/T^2$. The elementary
second-moment inequality gives
\[
 \Pr_0\{|M_T|\ge1/(2\sqrt T)\}\ge3/16.
\]
Indeed, Cauchy--Schwarz applied to that event gives
\[
 \Pr\{M_T^2\ge(\mathbb E M_T^2)/4\}
 \ge(3/4)^2\frac{(\mathbb E M_T^2)^2}{\mathbb E M_T^4}.
\]
Hence
\[
 H'(0)\ge
 T\frac1{32}\frac1{2\sqrt T}\frac3{16}
 =\frac{3\sqrt T}{1024}.
\]
For any exact sampler, expose only its distinct unknown input probes.
At $t=0$, the transcript score $Z$ has
$\mathbb E Z^2=\mathbb E Q$ and $H'(0)=\mathbb E[YZ]$, where
$Y\in\{-1,1\}$ is the output sign. Cauchy--Schwarz gives
$\mathbb E Q\ge H'(0)^2\ge9T/2^{20}\ge T/2^{17}$.
The average lower bound implies a worst-context lower bound. The
construction has $d+1=O(\log T)$ blocks and row norms at most four.
\end{proof}

The public constant coordinates make the denominator floor explicit.
The lower bound does not rely on normalization being undefined on any
input, on growing row norms, or on nonzero attention temperature. An
index built after reading the context can store all prefix means in this
example. It consequently evades this particular probe lower bound,
which is why cache preprocessing must be stated separately.


\begin{corollary}[A fresh-query obstruction survives a cache]
\label{cor:pn-fresh-query}
For every power-of-two width $n\ge2$, a pre-norm additive network with
one position, row norms at most four, stabilizer three, and
$O(\log n)$ blocks has a bounded binary output whose exact sampler
requires at least $n/2^{17}$ expected fresh-query coordinate probes in
the worst query. Any preprocessing independent of that query is allowed.
\end{corollary}
\paragraph{Proof idea.}
At one position, an affine attention value map averages the unknown
query coordinates. This places the same hard scalar recurrence inside a
single fresh query, so preprocessing independent of that query cannot
remove the information bound.

\begin{proof}
The query is $x\in\{-1,1\}^n$. Its first RMSNorm denominator is two.
At the sole position, attention has weight one. Give its value map the
matrix
\[
 V=2(J_n-I_n),\qquad (J_n)_{ij}=1/n.
\]
The first residual becomes $J_nx=M\boldsymbol1$, where
$M=n^{-1}\sum_i x_i$. Every row has norm $4(1-1/n)<4$, and the
coefficients are dyadic. Use a zero first tanh update, zero attention
updates thereafter, and identity tanh rows. All coordinates now obey
\[
 F_{r+1}(z)=F_r(z)+\tanh(F_r(z)/\sqrt{3+F_r(z)^2}).
\]
For $0\le y\le1$ this adds at least $y/4$, and for $y\ge1$ it adds at
least $1/4$. Choose $d$ exactly as in the preceding theorem with $n$ in
place of $T$, and use logits $\pm F_d(M)/(2d)$. The same moment and
finite-transcript argument gives $n/2^{17}$. An index built before $x$
arrives is part of the fixed preprocessed state in that argument.
\end{proof}

\subsubsection{Relative-floor lower bounds}

The next constructions keep a relative energy floor on the entire input
cube. They show that the polynomial dependence on the cumulative
residual radius in Theorem~\ref{thm:pn-relative} cannot be removed.

\begin{theorem}[A relative-floor depth lower bound, uniform in the step size]
\label{thm:pn-scaled-lower}
Fix a dyadic $s>1$ and put
\[
 \delta=s-1,\qquad C=\frac{s(s^2+1/2)}{2(s-1)}.
\]
There is a positive constant $c_s$, depending only on $s$,
with the following property.
For every power-of-two $N\ge2$, every dyadic $0<\alpha\le1$, and every
integer $D$ with $\alpha D\ge4$, there is a one-position, width-$2N$
pre-norm additive network with $D$ blocks, row norms at most $s$,
stabilizer $\epsilon=3$ in every normalization, and the same step
size $\alpha$ on every residual sublayer, whose exact binary output
requires
\begin{equation}
 \mathbb E Q\ge
 c_s
 \min\{(1+\alpha D)^{2\delta},N\}
 \label{eq:pn-scaled-lower}
\end{equation}
expected fresh-query coordinate probes in the worst input. Arbitrary
preprocessing independent of that query is allowed. On the entire input sign cube, every actual
normalization denominator squared satisfies
\[
 v_j\ge\frac1{12}(1+\alpha j)^2
\]
at the corresponding residual radius $R_j=1+\alpha j$.
Attention updates are zero, so the temperature is arbitrary. The
output uses one coordinate and a dyadic readout, and its logits lie
in $[-1,1]$ on every permitted input. The constant in
\eqref{eq:pn-scaled-lower} is given explicitly in the proof.
\end{theorem}
\paragraph{Proof idea.}
Separate feature coordinates from deterministic carriers that
control the normalization scale. Dividing by the residual radius gives a
scalar gain whose product grows polynomially. A reciprocal-square comparison
keeps that gain valid through nonlinear saturation, and the transcript score
turns it into a probe lower bound.

\begin{proof}
The network accepts every input in $\{-1,1\}^{2N}$. For its lower-bound
distribution, let the first $N$ coordinates be unknown signs
$x_1,\ldots,x_N$ and fix the remaining $N$ coordinates to $-1$. All query, key, and value maps are zero. Every feature tanh row
places weight $s/N$ on each of the first $N$ normalized coordinates,
with zero bias. Every carrier row has zero weights and bias $-1$.
These are legal dyadic parameters with the asserted row norms.

Let $M=N^{-1}\sum_i x_i$ and $t_0=\tanh1$. After $j$ blocks the feature
coordinates are
\[
 x_i-M+F_j(M),
\]
and all carrier coordinates are $C_j=-(1+\alpha jt_0)$. This follows
by induction, with $F_0(z)=z$ and
\begin{equation}
 F_{j+1}(z)=F_j(z)+\alpha
 \tanh\!\left(
 \frac{sF_j(z)}{
  \sqrt{3+[1-z^2+F_j(z)^2+C_j^2]/2}}
 \right).
 \label{eq:pn-scaled-meanrec}
\end{equation}
Indeed, the mean of the feature coordinates is $F_j(M)$, and their
mean square is $1-M^2+F_j(M)^2$. Every coordinate has magnitude at
most $R_j=1+\alpha j$. Since $t_0>1/2$,
$|C_j|\ge R_j/2$, so the denominator squared is at least $R_j^2/8$.
The attention residual has zero update and leaves these statements
unchanged. The network also has a relative floor on the entire sign cube.
For arbitrary initial carrier signs $e_i\in\{-1,1\}$, their values are
$e_i-\tau t_0$, with $\tau=\alpha j$. If $\tau\le2$, the stabilizer alone
gives $v_j\ge3\ge(1+\tau)^2/12$. If $\tau\ge2$, every carrier has
magnitude at least $\tau/2-1$, so
\[
 12v_j-(1+\tau)^2
 \ge12\left[3+\frac12(\tau/2-1)^2\right]-(1+\tau)^2
 =\frac12(\tau-8)^2+9>0.
\]
Thus the claimed floor $R_j^2/12$ holds on every sign input; the stronger
$R_j^2/8$ remains available on the lower-bound subdomain.

For $z\ge0$, all $F_j(z)$ are nonnegative and $F_j(z)\ge z$; the
recurrences are odd. Set $u_j=F_j/R_j$ and
$a_j=\alpha/R_{j+1}$. Their normalized update is
\[
 u_{j+1}=(1-a_j)u_j+
 a_j\tanh\!\left(\frac{s u_j}{\sqrt{b_j(z)+u_j^2/2}}\right),
\]
where
\[
 b_j(z)=\frac{3+(1-z^2)/2+C_j^2/2}{R_j^2}
 \le\frac{7}{2R_j^2}+\frac12.
\]
Let $j_0=\lceil2/\alpha\rceil$. Then $3\le R_{j_0}\le4$, hence
$u_{j_0}(z)\ge z/4$, and for every $j\ge j_0$ one has
$b_j(z)\le8/9<1$.

For $u\ge0$, the bound $\tanh y\ge y/\sqrt{1+y^2}$ gives
\[
 \tanh\!\left(\frac{s u}{\sqrt{b_j(z)+u^2/2}}\right)
 \ge\frac{s u}{\sqrt{1+(s^2+1/2)u^2}}.
\]
Put $g_j=1+\delta a_j$ and
$d_j=a_js(s^2+1/2)/g_j$. Convexity of
$v\mapsto(1+v)^{-1/2}$ gives
$u_{j+1}\ge g_ju_j/\sqrt{1+d_ju_j^2}$. Define
\[
 G=\prod_{j=j_0}^{D-1}g_j.
\]
For a positive $z$, write $G_j=\prod_{r=j_0}^{j-1}g_r$. The
reciprocal-square inequality gives
\[
 \frac{G_{j+1}^2}{u_{j+1}(z)^2}
 \le\frac{G_j^2}{u_j(z)^2}+d_jG_j^2.
\]
The initial reciprocal term is at most $16/z^2$. Moreover,
$d_j\le a_js(s^2+1/2)$ and $g_j^2-1\ge2\delta a_j$, so
\[
 \sum_{j=j_0}^{D-1}d_jG_j^2
 \le C\sum_{j=j_0}^{D-1}(g_j^2-1)G_j^2=C(G^2-1).
\]
Summing and inverting therefore gives
\begin{equation}
 u_D(z)\ge
 \frac{G(z/4)}{\sqrt{1+C(G^2-1)(z/4)^2}},\qquad z\ge0.
 \label{eq:pn-scaled-comparison}
\end{equation}
Continuity covers $z=0$. The calculation uses the pointwise comparison
at each step and does not assume that $b_j(z)$ is independent of $z$.

Choose the least power of two $L\ge R_D$. Then $R_D\le L<2R_D$.
The logits are plus and minus the first feature coordinate divided
by $L$. They lie in $[-1,1]$, and $1/L$ is dyadic. Conditional on
$M=z$, the selected sign $x_1$ has mean $z$. Put
\[
 k=1/L,\qquad A=(F_D(z)-z)/L\ge0.
\]
The conditional output-sign mean is
\[
 \overline g(z)
 =\frac{1+z}{2}\tanh(A+k)+
   \frac{1-z}{2}\tanh(A-k).
\]
At $A=0$, this equals $z\tanh k\ge zk/4$. Increasing $A$ from zero
to its actual value keeps both arguments in $[-1,1]$, since
$F_D(z)-z\le\alpha D=R_D-1$. Throughout that interval the derivative
of the displayed conditional mean with respect to $A$ is at least
$\operatorname{sech}^2(1)>1/4$. Consequently
\begin{equation}
 \overline g(z)\ge\frac{zk+A}{4}
 =\frac{F_D(z)}{4L}\ge\frac{u_D(z)}8,\qquad z\ge0.
 \label{eq:pn-scaled-coordinatehead}
\end{equation}
Oddness gives the signed counterpart for negative $z$.

Give the unknown query signs independent common mean $t$, and let
$H(t)$ be the output-sign mean. Differentiating the finite input sum
at zero and conditioning on $M$ gives
$H'(0)=N\mathbb E_0[M\overline g(M)]$. Apply
\eqref{eq:pn-scaled-comparison}--\eqref{eq:pn-scaled-coordinatehead}
and Jensen's inequality weighted by $M^2$. Since
$\mathbb E M^2=1/N$ and $\mathbb E M^4/\mathbb E M^2\le3/N$,
\[
 H'(0)\ge\frac{G}{32\sqrt{1+3CG^2/(16N)}}.
\]
The finite-transcript information inequality therefore gives
\[
 \mathbb E Q\ge
 \frac{G^2}{1024(1+3CG^2/(16N))}
 \ge\frac{\min\{G^2,N\}}{1024(1+3C/16)}.
\]
It holds on average over unbiased inputs, and hence for a worst input.

It remains to bound $G$ uniformly in $\alpha$. Monotonicity of
$x\mapsto(1+x)^{-1}$ and its square implies
\[
 \sum_{j=j_0}^{D-1}a_j
 \ge\log\frac{R_{D+1}}{R_{j_0+1}},\qquad
 \sum_{j=j_0}^{D-1}a_j^2
 \le\frac{\alpha}{R_{j_0}}\le\frac13.
\]
Also $R_{j_0+1}\le5$ and $R_{D+1}\ge1+\alpha D$.
The inequality $\log(1+x)\ge x-x^2/2$ gives
\[
 G\ge e^{-\delta^2/6}
       \left(\frac{1+\alpha D}{5}\right)^\delta.
\]
Since $e^{-\delta^2/3}5^{-2\delta}\le1$, substitution into the
minimum bound proves \eqref{eq:pn-scaled-lower} with
\[
 c_s=\frac{e^{-\delta^2/3}5^{-2\delta}}
              {1024(1+3C/16)}.
\]
The assumption
$\alpha D\ge4$ ensures $D\ge j_0$; all sums and comparisons above
are then valid.
\end{proof}


\begin{corollary}[A fixed final RMSNorm readout]
\label{cor:pn-scaled-finalnorm-lower}
In Theorem~\ref{thm:pn-scaled-lower}, replace the depth-dependent readout
by a final $\operatorname{RMSNorm}_3$. Use the fixed logits
$\pm\mathcal N_3(h_D)_1/4$. On the entire sign cube their absolute
values are at most $\sqrt{12}/4<1$. The lower bound
\eqref{eq:pn-scaled-lower} remains valid with $c_s/16$ in place
of $c_s$. Every normalization, including the final one,
has the same fixed stabilizer and the relative floor $R_D^2/12$.
\end{corollary}
\paragraph{Proof idea.}
The final normalization denominator stays between two fixed
multiples of the residual radius. Replacing the scaled readout by this
denominator therefore preserves the mean sensitivity up to a constant.

\begin{proof}
Conditional on $M=z$ in the lower-bound distribution, write
$L(z)=4\sqrt{v_D(z)}$. The full-domain floor gives $L(z)\ge R_D$.
The coordinate-radius bound gives $v_D(z)\le3+R_D^2\le4R_D^2$, hence
$L(z)\le8R_D$. The conditional shift calculation in
\eqref{eq:pn-scaled-coordinatehead} now gives
\[
 \overline g(z)\ge\frac{F_D(z)}{4L(z)}\ge\frac{u_D(z)}{32}
 \qquad(z\ge0).
\]
The denominator is even in $z$, so oddness remains valid. The score
lower bound loses a factor four, and its square loses a factor sixteen.
This gives the asserted constant. The full-domain logit bound is
$|h_{D,1}|/(4\sqrt{v_D})\le\sqrt{12}/4$, and the readout coefficient
$1/4$ is fixed and dyadic.
\end{proof}

\begin{corollary}[Padding to every width]
\label{cor:pn-scaled-padding}
There is a constant $c_s^{\rm pad}>0$, depending only on $s$, such
that for every integer $n\ge4$, the fixed final-normalization construction
can be made at width $n$, with all the assumptions of
Corollary~\ref{cor:pn-scaled-finalnorm-lower}, and with
\[
 \mathbb E Q\ge
 c_s^{\rm pad}
 \min\{(1+\alpha D)^{2\delta},n\}.
\]
An explicit value for the constant is given in the proof.
The relative floor $1/12$, fixed stabilizer $3$, and fixed readout
coefficient $1/4$ hold on the entire input sign cube.
\end{corollary}
\paragraph{Proof idea.}
Use a power-of-two feature group so its averaging weights remain
dyadic, and fill the remaining coordinates with carriers. Their positive
fraction preserves the energy floor, while the feature group still contains
a constant fraction of the unknown input.

\begin{proof}
Choose the largest power of two $N\le n/2$. Use $N$ feature
coordinates and $n-N$ carrier coordinates. Put $\eta=N/n$, so
$1/4<\eta\le1/2$. Each feature row still has the dyadic weights
$s/N$. On the lower-bound subdomain the squared denominator is
\[
 3+\eta(1-z^2+F_j(z)^2)+(1-\eta)C_j^2.
\]
The carrier fraction is at least $1/2$, so the full-cube floor proof
above is unchanged. Use $j_0=\lceil3/\alpha\rceil$. Then
$4\le R_{j_0}\le5$, $u_{j_0}\ge z/5$, and the normalized nonsignal
part of the denominator satisfies
\[
 b_j(z)\le\frac{7}{2R_j^2}+\frac34\le\frac{31}{32}<1.
\]
Because $\eta\le1/2$, the same Jensen comparison applies with the
same $C$, replacing $z/4$ by $z/5$. The fixed final-normalization
head still gives $\overline g(z)\ge u_D(z)/32$. Thus the score and
query inequalities become
\[
 H'(0)\ge\frac{G}{160\sqrt{1+3CG^2/(25N)}},\qquad
 \mathbb E Q\ge\frac{\min\{G^2,N\}}{25600(1+3C/25)}.
\]
Now $\sum a_j^2\le1/4$ and $R_{j_0+1}\le6$, whence
$G\ge e^{-\delta^2/8}((1+\alpha D)/6)^\delta$.
Finally, $N>n/4$ implies
$\min\{x,N\}\ge\min\{x,n\}/4$ for $x\ge0$. These estimates
give the explicit constant
\[
 c_s^{\rm pad}=\frac{e^{-\delta^2/4}6^{-2\delta}}
                         {102400(1+3C/25)}.
\]
Since $\alpha D\ge4$,
$D\ge j_0$ as required.
\end{proof}

\begin{corollary}[Prescribed bounded step schedules]
\label{cor:pn-scheduled-lower}
The arbitrary-width construction in Corollary~\ref{cor:pn-scaled-padding}
also permits prescribed dyadic block steps $a_j\in[0,1]$ for $0\le j<D$,
used on both residual sublayers of their respective blocks. Put
$\tau_j=\sum_{k<j}a_k$, and suppose $\tau_D\ge4$.
For a positive constant $c_s^{\rm sched}$ depending only on $s$,
its fixed final readout requires
\[
 \mathbb E Q\ge
 c_s^{\rm sched}
 \min\{(1+\tau_D)^{2\delta},n\}.
\]
The constant is specified in the proof.
All full-cube floors hold at radius $r_j=1+\tau_j$ with the same
constant $1/12$.
\end{corollary}
\paragraph{Proof idea.}
The carrier displacement depends on the cumulative step size.
Reciprocal increments telescope even for an irregular schedule, and their
squares have a bounded sum. These two facts control the product of normalized
gains without a regularity assumption on the steps.

\begin{proof}
Replace every occurrence of $\alpha j$ in the carrier and feature
recurrences by $\tau_j$, and every increment $\alpha$ by $a_j$.
Let $j_0$ be the first index with $\tau_{j_0}\ge3$. Since $a_j\le1$,
$4\le r_{j_0}\le5$, so the same initial value $u_{j_0}\ge z/5$
and nonsignal denominator bound apply. The normalized branch size is
$q_j=a_j/r_{j+1}$. Let $G=\prod_{j=j_0}^{D-1}(1+\delta q_j)$.
The reciprocal-square comparison and the score lower bound are
unchanged.

For completeness, the needed product estimate follows without any
regularity assumption on the schedule. Since $r_{j+1}-r_j=a_j$,
\[
 q_j^2\le\frac{a_j^2}{r_jr_{j+1}}
 \le\frac1{r_j}-\frac1{r_{j+1}}.
\]
Also, integrating $1/r-1/r_{j+1}$ over $[r_j,r_{j+1}]$ gives
\[
 0\le\log\frac{r_{j+1}}{r_j}-q_j
 \le\frac{a_j^2}{r_jr_{j+1}}.
\]
Therefore $\sum q_j^2\le1/4$ and
$\sum q_j\ge\log(r_D/r_{j_0})-1/4$. The inequality
$\log(1+x)\ge x-x^2/2$ yields
\[
 G\ge e^{-\delta/4-\delta^2/8}
       \left(\frac{1+\tau_D}{5}\right)^\delta.
\]
Substitution into the score lower bound proves the claim with
\[
 c_s^{\rm sched}=\frac{e^{-\delta/2-\delta^2/4}5^{-2\delta}}
                           {102400(1+3C/25)}.
\]
Zero steps
leave the state unchanged and cause no exception in these identities.
\end{proof}

\begin{corollary}[An effective-depth threshold for the certified pre-norm class]
\label{cor:pn-effective-threshold}
Fix $s>1$, a nonnegative temperature $\beta$, a rational
$0<\gamma\le1/12$, an output row-norm bound $u\ge1/4$, and an
output bias bound $c_{\rm out}\ge0$, all independent of width and depth.
Consider widths tending to infinity, polynomially
bounded position counts and depths, and encoding/address length
$B=\log^{O(1)}n$. Use a common dyadic residual step $\alpha\in[0,1]$.
For each attention update take the valid radius
$U_j=G\max_i\sum_k|V_{j,ik}|$, and for a tanh update take $U_j=1$,
omitting updates certified to be zero. Let
$R_j=1+\alpha\sum_{k\le j}U_k$. Restrict to networks whose pre-norm
and final-normalization denominators satisfy the relative certificate
$v_j\ge\gamma R_{j-1}^2$ of \eqref{eq:pn-relative-floor}, with
$j=2D+1$ for the final normalization. The affine output logits lie in $[-1,1]$;
their row norms are at most $u$ and their biases have magnitude
at most $c_{\rm out}$.

This class has a uniform polylogarithmic worst-case exact-sampling
bound if and only if
\[
 \alpha D=\log^{O(1)}n.
\]
The equivalence concerns asymptotic families. It determines the scale
of effective depth; it does not identify the optimal polynomial
exponent.
\end{corollary}
\paragraph{Proof idea.}
The composition theorem gives the upper bound from the cumulative
residual radius. For the converse, one fixed gain above one yields a
positive lower exponent on every width, so a radius larger than every
polylogarithm cannot have a uniform polylogarithmic sampling bound.

\begin{proof}
The positive direction follows from
Corollary~\ref{cor:pn-stepsizes}, since $U_j\le sG$ at attention and
$U_j\le1$ at tanh, while all local factory constants are fixed.
The final-normalization interface has fixed radius and cost.

For necessity, choose a fixed dyadic $\sigma$ with $1<\sigma\le s$
(if $s$ is dyadic, take $\sigma=s$). For every width $n\ge4$,
use Corollary~\ref{cor:pn-scaled-padding} with gain $\sigma$. Its attention update
radii are zero and its tanh update radii are one, so the certificate
radii are exactly $1+\alpha j$. Its full-domain floor $1/12$ meets
any stated $\gamma\le1/12$. Its fixed coefficient $1/4$ gives bounded
affine logits. The lower bound is a constant depending only on $s$
times $\min\{(1+\alpha D)^{2(\sigma-1)},n\}$. If $\alpha D$ is not bounded
by a fixed power of $\log n$, this lower bound is not uniformly
polylogarithmic. The regime $\alpha D<4$ cannot cause that failure,
so the lower theorem's cutoff has no effect on the equivalence.
\end{proof}
